\subsection*{唯一性条件下的Euler折线}
Peano 定理利用紧性从 Euler 折线中抽取收敛子列。若另有 Lipschitz 条件，则第~\ref{sec:picard-contraction}~节的残差估计还能控制整个近似序列。
必要时缩小共同存在区间，仍记其半长为 $h$，并使 $Lh<1$。在该区间的连续函数闭集上，令
\[
 (Au)(t)=x_0+\int_{t_0}^tf(s,u(s))\,ds.
\]
记真解为 $x$，则 $Ax=x$，且对任一留在该闭集内的近似曲线 $u$，
\[
 \|u-x\|_\infty\leq\|u-Au\|_\infty/(1-Lh).
\]
在 Peano 定理的证明中已经得到 Euler 折线 $u_n$ 满足
\[
 u_n=Au_n+\sigma_n,\qquad\|\sigma_n\|_\infty\longrightarrow0.
\]
代入残差估计，便有
\[
 \|u_n-x\|_\infty\leq\frac{\|\sigma_n\|_\infty}{1-Lh}\longrightarrow0.
\]
因此在这个共同小区间上，整个 Euler 折线序列均一致收敛到解，而不只是某个子序列。
